% ---------------------------------------------------------------------
%  Appendix A : Glossary
%  A single alphabetised glossary synthesised from the per-subsystem
%  briefs.  Every term that appears anywhere in the manual is defined
%  here from scratch, for a reader who knows the physics but not the
%  specific tooling (SageMath, nauty, sympy, numba, ...).
% ---------------------------------------------------------------------
\chapter{Glossary}\label{ch:glossary}

This appendix collects, in one alphabetised list, every piece of jargon the
manual leans on. It is deliberately written so that you can land here from a
cross-reference in any chapter and walk away understanding the term \emph{cold},
without having read the chapter it came from. Three kinds of vocabulary are mixed
together on purpose, because in practice you meet them mixed together:

\begin{itemize}
  \item \textbf{Physics and field theory} --- the \msrjd{} action, propagators,
        cumulants, loop order, symmetry factors, Symanzik polynomials. A reader
        who knows stochastic field theory will recognise these but may not know
        exactly how \Daedalus{} names or splits them.
  \item \textbf{Mathematical machinery} --- Schwinger parameters, divided
        differences, the Hermite--Genocchi formula, causal posets, chain
        simplices, Cholesky and Lyapunov solves. These are the tools the loop
        integrals actually reduce to.
  \item \textbf{Software tooling} --- SageMath, \code{nauty}, \code{sympy},
        \code{numba}, \code{scipy}, \code{ipywidgets}, and the dozens of small
        library calls (\code{fast\_callable}, \code{lambdify},
        \code{solve\_continuous\_lyapunov}, \code{importlib.util}, \ldots) that
        do the work. This is the layer the manual promises to over-explain, so it
        gets the most attention here.
\end{itemize}

\noindent
A handful of the most load-bearing entries (the \msrjd{} action, the symmetry
factor, the Symanzik polynomials, Phase~J) are given a boxed \emph{Definition}
rather than a one-line gloss, because you will meet them on nearly every page.
Where a term names a concrete function, class, or attribute in the codebase, the
defining file is cited so you can read the implementation. Identifiers are set in
\code{monospace}; the field-theory ``response'' field is written
$\tilde\varphi$ in the maths and, by the engine's own convention, with a
\code{t}- or \code{d}-prefix in code (\code{xt}, \code{dn}).

The entries are grouped by first letter. Symbols and Greek-letter terms are
filed under the letter of their romanised name (so $\Sym(\Gamma)$ appears under
\textbf{S}, for ``symmetry factor''; $\kappa^{(n)}$ under \textbf{C}, for
``cumulant'').

% =====================================================================
\section*{A}
% =====================================================================
\begin{description}[leftmargin=0pt,style=nextline]

  \item[\code{abs\_symbolic}] Sage's symbolic absolute-value function. It keeps an
        expression such as $|\tau|$ \emph{unevaluated} as a symbol, rather than
        forcing a numeric branch. The colored-noise kernels need this because the
        Lorentzian $\ee^{-|\tau|/\tau_c}$ must survive as an even function of the
        time lag all the way into the Fourier transform.

  \item[Action $S$] The functional in the exponent of the \msrjd{} path integral,
        $\ee^{-S[\varphi,\tilde\varphi]}$. It is the central object the user
        authors --- as a text string passed to \code{.set\_action\_text(...)}. Its
        stationary point is the classical (mean-field) solution; its fluctuations
        about that point generate every correlation function. The bilinear part is
        the free propagator; the cubic-and-higher part supplies the interaction
        vertices.

  \item[\code{add\_vararg}] Sage's internal $n$-ary addition operator. It is
        \emph{not} Python's \code{operator.add} (which is binary); when the
        field-theory core walks a symbolic expression tree it detects a sum by
        matching this operator by name.

  \item[Adjugate $\mathrm{adj}(K)$] The transpose of the cofactor matrix of a
        square matrix $K$. It gives the inverse through Cramer's rule,
        $K^{-1}=\mathrm{adj}(K)/\det(K)$, which is how the propagator
        $G=K^{-1}$ is formed symbolically before any numbers are substituted.

  \item[Allen--Cahn / reaction--diffusion] The prototypical first-tier spatial
        model: a field with a relaxational mass $\mu$ and a diffusion term
        $D\nabla^{2}$. Its inverse propagator has the form $A + B k^{2} + \ii\omega$.

  \item[\code{analytic-eligible}] A property of a $\delta$-edge subset (temporal)
        or a diagram (spatial): the loop integral has a closed form and needs no
        \code{scipy.nquad} fallback. A subset is analytic-eligible exactly when no
        smooth-kernel noise vertex (a \code{NoiseSourceType} with a non-rational
        kernel) is present; otherwise the integrand is non-rational and the
        adaptive numerical path is forced.

  \item[\code{ast} (Abstract Syntax Tree)] Python's standard-library module for
        parsing source text into a tree of syntax nodes \emph{without running it}.
        The UI and the theory loader use it to validate an action string --- they
        walk the parsed tree to check that only allowed names appear, never
        executing the code.

  \item[\code{ast.literal\_eval}] A safe evaluator that turns a string into a
        Python value \emph{only} if the string is a literal (number, string, list,
        dict, tuple, bool, or \code{None}). It never executes function calls or
        arbitrary code, which is why it is used for the random-seed box and similar
        user inputs.

  \item[Automorphism group $\Aut(\Gamma)$] The group of relabellings of a graph
        $\Gamma$ that map it onto itself edge-for-edge (preserving any colours).
        Two variants matter for symmetry factors: \code{Aut\_fixed\_ext} pins the
        external leaves in place, while \code{Aut\_free\_ext} (also written
        \code{Aut\_free}) lets same-field external leaves permute. Their ratio is
        the external-Wick compensation.

  \item[Auxiliary field (aux field)] An extra field injected by Markovian
        embedding to represent colored noise: the Ornstein--Uhlenbeck field
        $\xi$ (\code{xi}) together with its \msrjd{} response partner
        $\tilde\xi$ (\code{xit}). Adding it enlarges the state space so the dynamics
        become local in time.

\end{description}

% =====================================================================
\section*{B}
% =====================================================================
\begin{description}[leftmargin=0pt,style=nextline]

  \item[Betti number (first) / loop number] For a connected graph,
        $|E|-|V|+1$ --- the number of independent cycles, equivalently the loop
        order $\ell$ of the corresponding diagram.

  \item[Bigrade $(n_{\tilde{}},\,n_{\mathrm{phys}})$] The grading of an action
        monomial by $(\text{number of response factors},\ \text{number of
        physical factors})$; the total degree is their sum. This single pair is
        the primary classification axis of the expanded action: the saddle sector
        is $(0,0),(1,0),(0,1)$; the free propagator kernel is $(1,1)$; the white
        noise lives at $(2,0)$; any monomial of total degree $\ge 3$ is an
        interaction vertex. In the \code{FieldTheory} object the attribute
        \code{\_by\_tp} maps each bigrade to the sum of its monomials.

  \item[Bilinear vs vertex generator] In the operator-IR authoring path, a
        distinction between IR generators that fold into the propagator (total
        degree $\le 2$, the bilinear part) and those that carry per-leg momentum
        form factors forward into the loop integrator (total degree $\ge 3$, the
        vertices).

  \item[\code{bliss}] A graph-automorphism engine, an alternative to \code{nauty};
        Sage may call either internally for \code{automorphism\_group}. (See
        \code{nauty}.)

  \item[Bubble] A one-loop self-energy insertion with two internal vertices joined
        by two edges; its loop momentum couples to the external momentum $q$. In a
        multigraph a bubble is a double edge between a vertex pair.

\end{description}

% =====================================================================
\section*{C}
% =====================================================================
\begin{description}[leftmargin=0pt,style=nextline]

  \item[Canonical form / canonical label] A deterministic relabelling of a graph's
        vertices such that two graphs are isomorphic if and only if their canonical
        forms are \emph{identical}. This converts the expensive ``are these the
        same diagram?'' question into a cheap string comparison. Sage exposes it as
        \code{DiGraph.canonical\_label}; the routine and the original$\to$canonical
        id map (the \emph{certificate}, \code{cert}) come from \code{nauty}.

  \item[Catastrophic cancellation] The loss of significant digits that occurs when
        you subtract two nearly-equal large floating-point quantities. It is the
        reason the symbolic adjugate residues are unreliable close to a pole, and
        why the propagator code switches to an exact rational field there.

  \item[Causal chamber / chamber] One ordering region of the internal-vertex-time
        integral, corresponding to a single linear extension of the retarded
        partial order. Within a chamber every time-difference sign is fixed, so the
        integrand is smooth and the integral is closed-form. The full loop integral
        is the (disjoint) sum over all chambers. (Spatial pipeline; the temporal
        analogue is the \emph{polytope}/\emph{chain-simplex} decomposition.)

  \item[Causal filtering] Discarding the typed diagrams that violate the retarded
        \msrjd{} propagator structure --- those that would require an effect before
        its cause. Implemented in \code{filter\_causal}
        (\file{msrjd/diagrams/causality.py:126}). A causal diagram is structurally
        a directed acyclic graph (DAG).

  \item[Causal poset / causal polytope] For three or more surviving integration
        times, the set of retardation constraints $\Delta t_e > 0$ read as a
        partial order ($s_u < s_v$). The convex region it carves out is the
        \emph{causal polytope}, the integration domain; integrating over its
        linear extensions (nested \emph{chain simplices}) evaluates the polytope
        integral.

  \item[\code{CDF} (Complex Double Field)] Sage's type for machine-precision
        (hardware \code{double}) complex numbers. The pipeline drops into \code{CDF}
        for the fast numeric phases --- finding poles, evaluating residues --- after
        the exact symbolic work is done.

  \item[\code{C} edge / correlation line] An edge carrying a
        $\langle\varphi\varphi\rangle$ correlation. In the heat-kernel
        representation it is two retarded ($\code{R}$) segments glued at a
        two-point noise source, parameterised by a single Schwinger time
        $\sigma\in[0,\infty)$. Contrast the \code{R} edge.

  \item[Certificate (\code{cert})] \code{nauty}'s original-to-canonical vertex id
        map produced alongside the canonical label; equal canonical forms (and
        their certificates) certify two diagrams as isomorphic. (See canonical
        form.)

  \item[CGF / CGF term / CGF row] A Cumulant-Generating-Functional term: a single
        declared contribution to the noise's cumulants, given as a tuple
        \code{(coefficient, kernel, order, response\_legs)} via
        \code{.declare\_cgf\_term(...)}. Order $2$ is the noise covariance; orders
        $\ge 3$ are non-Gaussian noise. ``CGF'' is also the internal name of the
        UI's Noise tab.

  \item[Chain simplex] A nested-inequality region
        $L \le s_1 \le \dots \le s_N \le U$. The exponential integral over a chain
        simplex has a closed form with $2^{N}$ terms; the temporal $m\ge 3$
        integrator sums over the chain simplices that tile the causal polytope.

  \item[Characteristic polynomial / \code{D\_omega}] $\det(K_{\mathrm{ft}})$, the
        determinant of the Fourier-image kernel. The roots of its numerator are the
        system's poles $\omega_k$.

  \item[Cholesky factorization] The decomposition $C=LL^{\mathsf T}$ of a symmetric
        positive-semidefinite matrix. Multiplying a vector of independent unit
        Gaussians by $L$ yields correlated Gaussians with covariance $C$ --- the
        trick a simulator uses to draw correlated noise (the explicit $2\times2$
        form gives $\eta_y=\rho u + \sqrt{1-\rho^{2}}\,v$).

  \item[Cofactor expansion] Computing a determinant or adjugate by expanding along
        minors. The dense NumPy path forms each minor with \code{np.delete}.

  \item[Coloured incidence digraph] The bipartite graph built from a diagram so
        that an isomorphism engine can see edge types: vertices and \emph{edges}
        both become nodes, each coloured by its field / vertex / propagator type.
        This is the object fed to \code{nauty} for canonical labelling and
        automorphisms. (Also ``incidence digraph (coloured)''.)

  \item[Colored noise] Noise with a finite correlation time $\tau_c$; its
        covariance is a smooth function of the time lag (here the Lorentzian
        $c\,\ee^{-|\tau|/\tau_c}$) rather than a $\delta$-function. Contrast white
        noise.

  \item[Completeness bound $(k+3\ell-j-1)$] The proven upper bound on the number of
        degree-2 vertices that a valid topology can have; respecting it guarantees
        the tree search misses no diagram. It sets \code{v2\_max}. (An older
        $-\lfloor j/3\rfloor$ form of the lemma was false at $\ell\ge 3$; the
        bound above is the corrected, proven one.)

  \item[Composite vs per-leg vertex] Two modes of a derivative interaction vertex.
        A \code{'composite'} vertex applies the operator to the whole
        $\varphi^{n}$ composite, giving a response-leg momentum factor (e.g.
        $\nabla^{2}(\varphi^{2})$ in Model~B, $\partial_x(\varphi^{2})$ in
        Burgers). A \code{'perleg'} vertex applies the operator to each physical
        leg separately (e.g. $(\partial_x\varphi)^{2}$ in KPZ, giving
        $\prod \ii\,p_{\mathrm{leg}}$). The two produce different form factors.

  \item[Convolution \code{Conv(g, n)}] The formal time-convolution operator
        $\int g(t-s)\,n(s)\,\dd s$. It is asymmetric --- only its second argument
        carries time content --- and each downstream consumer reduces it
        differently. Reduced symbolically via Sage's \code{substitute\_function}.

  \item[Convolution theorem] $(g\star n)\widehat{\ }(\omega)=\hat g(\omega)\,\hat
        n(\omega)$: convolution in time is multiplication in frequency. This is why
        a kernel can be left as an abstract symbol and only resolved to its
        frequency image $\hat g(\omega)$ at Fourier-transform time.

  \item[Crank--Nicolson / semi-implicit] A trapezoidal-rule implicit time-stepping
        scheme. In the coupled reaction--diffusion simulator it delivers the
        $O(\dd t^{2})$-accurate stationary covariance through a Pad\'e$(1,1)$
        propagator.

  \item[Cumulant / connected correlation function] A connected $k$-point function
        $\kappa_k=\langle\varphi\cdots\varphi\rangle_c$ --- the sum of all
        \emph{connected} Feynman diagrams with $k$ external legs. This is the
        native output of \code{compute\_cumulants}. The second cumulant of the
        noise is its covariance $\langle\xi\xi\rangle$.

  \item[$\kappa^{(n)}$ (noise cumulant)] The $n$-th cumulant of a noise process:
        $\kappa^{(2)}$ is Gaussian (the covariance), $\kappa^{(3)},\kappa^{(4)},\dots$
        are higher (non-Gaussian) corrections. Correlated noise contributes a
        cumulant series $-W_m[\tilde m]$ to the action.

  \item[\code{\_cumulant\_kernels}] A namespace dict of noise-cumulant kernel
        callables. It is a \emph{runtime} side effect of
        \code{\_build\_cumulant\_action} and is deliberately not pickled --- it is
        rebuilt when a cache is loaded.

  \item[\code{CyclotomicField(4)} $=\mathbb{Q}(\ii)$] The exact number field
        obtained by adjoining the imaginary unit to the rationals. Working in
        $\mathbb{Q}(\ii)$ (rather than floating point) gives reliable polynomial
        GCDs, which is why it is used for the canonical, exact propagator
        inversion.

  \item[\code{cysignals} (\code{alarm} / \code{SignalError} / \code{AlarmInterrupt})]
        Cython tools that let Python time-bound and recover from long-running
        \emph{native} routines (Singular, Pynac). They are how the propagator code
        puts a wall-clock budget on a symbolic simplification that might otherwise
        hang. (See also \code{SIGALRM}.)

\end{description}

% =====================================================================
\section*{D}
% =====================================================================
\begin{description}[leftmargin=0pt,style=nextline]

  \item[DAE (differential-algebraic equation)] A system mixing differential
        equations (those carrying the time-derivative \code{Dt}) with algebraic
        constraints (those without). The mean-field saddle is the solution of the
        DAE obtained from \code{.equation(lhs, rhs, population)} after substituting
        \code{Dt} $\to 0$.

  \item[DAE saddle solve / fixed-point index] A multi-root Newton solver for the
        algebraic mean-field system. When a theory has several stable saddles (e.g.
        a double well), \code{fixed\_point\_index} selects which sorted root the
        loop expansion is built around.

  \item[DAG (Directed Acyclic Graph)] A directed graph with no directed cycle. A
        causal (retarded) diagram is structurally a DAG; this is the graph-theoretic
        face of causality.

  \item[\code{D\_delta} ($\delta$-coefficient / instantaneous part)] The
        $\omega\to\infty$ limit of the Fourier propagator $\hat G(\omega)$ --- the
        matrix weight of an instantaneous $\delta(t)$ term in the time-domain
        propagator. It is the ``immediate response'' piece.

  \item[DC gain $\hat g(0)$] A kernel's integral $\int g(t)\,\dd t$, i.e. its
        Fourier transform at zero frequency. For a unit-integral (normalised)
        kernel, $\hat g(0)=1$.

  \item[Dedup multiplicity] The size of a diagram's isomorphism-equivalence class.
        Under the framework's ``Path~A'' convention this number is \emph{diagnostic
        only}: it is reported but must not be multiplied back into the diagram's
        weight, because the symmetry factor already accounts for it.

  \item[Degree] The number of edge-ends meeting at a vertex; a multi-edge counts
        with its multiplicity.

  \item[$\delta$-edge / smooth-edge subset] The $2^{|E|}$-fold expansion of a
        product of retarded propagators, choosing on each edge whether to take its
        instantaneous ($\delta$) part or its smooth (exponential) part. Each choice
        is a ``subset''.

  \item[$\delta$-elimination] Using a $\delta$-edge --- which imposes $\Delta t=0$,
        a linear equation among the times --- to remove one integration variable.
        Solved symbolically with \code{sage\_solve}.

  \item[Derivative symbol] A renamed partial-derivative-at-saddle atom. Sage's
        formal differentiation produces clumsy ``derivative-of-function-at-0''
        objects; the field-theory core renames them to clean symbols such as
        \code{phi1\_2} ($=\varphi'_2(0)$) or \code{phi21\_1}
        ($=\partial^{2}_{x_1}\partial_{x_2}\varphi_1(0)$).

  \item[Derived generator] In the operator IR, a fresh polynomial-ring generator
        introduced to stand for $\mathrm{Op}(\delta\varphi)$ --- the trick
        $u=\delta\varphi,\ v=\nabla^{2}\delta\varphi$ --- so that multivariate
        Taylor expansion treats the derivative quantity as if it were an ordinary
        field.

  \item[\code{diagram\_signature}] The canonical form of a typed diagram's coloured
        incidence digraph; a \emph{complete} isomorphism invariant. Two diagrams are
        duplicates if and only if their signatures match
        (\file{msrjd/diagrams/symmetry.py:467}).

  \item[Diffusion matrix $\mathcal{D}$] The $k^{2}$ part of the linearised inverse
        propagator of a multi-field spatial theory; it encodes spatial spreading.
        $\mathcal{D}\propto I$ means scalar (equal) diffusion across fields.

  \item[\code{DiGraph} / \code{Graph}] Sage's directed and undirected graph classes.
        Diagram topologies are built and manipulated as these objects; their
        \code{canonical\_label} and \code{automorphism\_group} methods are the
        bridge to \code{nauty}.

  \item[\code{dirac\_delta(tau)}] Sage's symbolic Dirac delta. It is the white-noise
        kernel --- the time-domain covariance of $\delta$-correlated noise.

  \item[Dirty flag / dirty guard] In the UI, the boolean \code{\_dirty} that tracks
        unsaved edits. Reset and Load are gated behind a confirmation checkbox while
        the form is dirty, so you cannot silently discard changes.

  \item[{Divided difference $f[x_0,\dots,x_n]$}] A classical finite-difference object.
        In the coupled spatial path the $n$-fold time-ordered (simplex) integral
        equals $(-1)^{n}$ times the $n$-th divided difference of
        $f(z)=\ee^{-tz}$ --- which is what makes the Dyson dressing computable in
        closed form.

  \item[Downgrade (filter)] Keeping only the bigrades of total degree
        $\le \text{target}$ from a higher-order cached expansion. It is exact by the
        superset theorem, and is how a cache built at one Taylor order serves a
        cheaper run.

  \item[Drift $V$] The $k^{1}$ (advection) coefficient of a spatial inverse
        propagator, $\ii v$ for a velocity $v$; it is zero for even / Laplacian-only
        kernels.

  \item[\code{'dn'} / \code{'dv'} / \code{'dm'}] Estimator field-type labels in the
        simulators: spike-train fluctuation (factorial-corrected),
        voltage fluctuation (smooth, linearly centred), and external-drive spike
        fluctuation (shot-noise-like, treated as \code{'dn'}).

  \item[Duhamel convolution] The time-ordered nested integral that expresses a
        perturbed time-evolution as a series of insertions sandwiched between
        unperturbed evolutions. The Dyson--Duhamel series is built from these.

  \item[\code{Dt}] The formal time-derivative operator symbol, usable directly in
        action and equation text. It is substituted to $0$ at the stationary saddle
        and to the eigenvalue $\sigma$ in linear-stability analysis.

  \item[Dyson dressing / Dyson--Duhamel series] Resumming a self-energy (or residual
        diffusion) series into the propagator, order by order, to a finite
        \code{dyson\_order}. Each order is an $n$-fold Duhamel convolution. Used for
        coupled fields with \emph{unequal} diffusivities, where the residual
        diffusion $\hat{\mathcal{D}}$ is the perturbation. (Also ``Dyson--Duhamel
        truncation''.)

\end{description}

% =====================================================================
\section*{E}
% =====================================================================
\begin{description}[leftmargin=0pt,style=nextline]

  \item[\code{erf}-split closed form (Rescue A)] The antiderivative of
        $s^{-1/2}\ee^{-\beta/s-\alpha s}$, written as a sum of two error functions
        with $\ee^{\pm 2\sqrt{\alpha\beta}}$ weights. It is the spatial heat-kernel
        integrator's closed form for a Schwinger integral that would otherwise need
        quadrature.

  \item[ETD1 (exponential time-differencing, order 1)] A time-stepping scheme for
        stiff equations $\partial_t\varphi=L\varphi+N(\varphi)$: it treats the
        linear operator $L$ \emph{exactly} (mode by mode in Fourier space) and adds
        the nonlinear part $N$ through the factor $(1-\ee^{-\omega\,\dd t})/\omega$.

  \item[Euler--Maruyama] The first-order numerical scheme for stochastic
        differential equations: $x \mathrel{+}= \dd t\cdot\text{drift} +
        \sqrt{2D\,\dd t}\,\mathcal{N}(0,1)$. The $\sqrt{\dd t}$ scaling of the noise
        term is what distinguishes it from deterministic Euler stepping.

  \item[\code{eval}] Python's built-in expression evaluator. The DAE solver
        \code{eval}s equation text inside a deliberately sandboxed namespace (only
        the field and parameter symbols are exposed).

  \item[External fields / legs] The $k$ field operators whose correlator is being
        computed. In code, \code{external\_fields} is a length-$k$ list of
        \code{(field\_name, slot)} tuples; \code{k} is just its length (the
        correlator order).

  \item[External-Wick compensation \code{comp}] The orbit--stabilizer divisor
        $|\Aut(\Gamma,\text{free})| / |\Aut(\Gamma,\text{fixed})|$ that removes the
        over-counting introduced when the integrator sums over permutations of
        same-field external leaves. (See orbit--stabilizer theorem.)

\end{description}

% =====================================================================
\section*{F}
% =====================================================================
\begin{description}[leftmargin=0pt,style=nextline]

  \item[Factorial moment / falling factorial] The product
        $n^{(m)}=n(n-1)\cdots(n-m+1)$. For a Poisson count $n$,
        $\langle n^{(m)}\rangle=\lambda^{m}$. Replacing $n^{m}$ by $n^{(m)}$ at
        coincident time bins strips the self-spike (shot-noise) $\delta$-terms from
        a point-process correlator, isolating the smooth off-diagonal cumulant the
        diagrams predict.

  \item[Fan triangulation] Splitting a convex polygon into triangles all sharing
        vertex $0$. It is how the temporal $m=2$ polytope integral is decomposed.

  \item[\code{fast\_callable}] Sage's compiler that turns a symbolic expression into
        a fast Python closure (Sage's analogue of \code{lambdify}). It is the
        slow-path integrand compiler --- ``slow'' only relative to the JIT-compiled
        inner loops, not to symbolic evaluation.

  \item[\code{FieldTheory} / \code{expand()}] The Phase-1 engine that
        Taylor-expands the symbolic action about the saddle and classifies every
        term by bigrade. \code{ensure\_taylor\_order} re-invokes \code{expand()}
        only when a cached expansion's order is too low for the requested diagrams.

  \item[\code{\_FieldScalar}] A size-1 wrapper class that supports both bare
        arithmetic and indexed access, so a user can write \code{xt*x} and
        \code{xt[0]*x[0]} interchangeably for a single-component field.

  \item[Fixed-point index] See \emph{DAE saddle solve}. The integer that picks which
        sorted mean-field root the diagrammatics expand around.

  \item[Fluctuation field] The internal field $\delta\varphi=\varphi-\varphi^{*}$
        --- the actual fluctuating quantity about the mean-field background, and the
        variable the action is Taylor-expanded in. Natural code names carry a
        \code{d}-prefix: \code{dn}, \code{dv}, \code{dx}.

  \item[Fluctuation--regression (Onsager regression) theorem] The statement that an
        equilibrium fluctuation relaxes by the same law as a macroscopic
        perturbation: $C(\tau)=\ee^{-A\tau}\Sigma$. It connects the stationary
        covariance $\Sigma$ to the time-dependent correlator.

  \item[Fork vs thread parallelism] Two ways to run work in parallel. \emph{Fork}
        clones the parent process so workers inherit its memory with no input
        pickling --- fast, but it crashes a macOS Jupyter kernel once Cocoa/BLAS
        have initialised. \emph{Threads} share one process's memory and are safe
        because heavy NumPy operations release the GIL. The temporal stages fork
        (behind a notebook-detecting guard); the spatial stages thread. \emph{Spawn}
        is a third method that pickles its arguments --- it would fail here on the
        unpicklable integrand closures.

  \item[Form factor] The polynomial in the momenta that a \emph{derivative}
        interaction vertex deposits on the loop integrand:
        $\nabla^{2}\to-|k|^{2}$, $\partial_x\to \ii k$,
        $\partial_t\to-\ii\omega$. Per-leg vs composite vertices give different
        form factors; in the coupled path it is written
        $\mathcal{R}(q,\ell)=\prod_{\text{vertices}}\mathfrak{F}(v)$.

  \item[Form-factor signature] \code{vertex\_form\_factor\_signature}: a SHA-256
        fingerprint of the operator-IR vertex table
        (\file{pipeline/\_expand\_cache.py:208}). It guards the expand cache against
        a slug that under-determines the per-vertex momentum form factors --- two
        models with the same name but different derivative vertices get different
        signatures.

  \item[Formal / undefined function] A Sage symbolic function such as
        \code{function('phi')(v)} left \emph{unapplied} --- with no concrete body
        --- so that \code{taylor()} can differentiate it abstractly. Its
        derivative-at-0 atoms are afterwards renamed to clean derivative symbols.
        (Also ``formal function''.)

  \item[Free action / bilinear kernel] The Gaussian $(1,1)$ sector of the action:
        exactly one response leg and one physical leg. \code{ft.free\_action()}
        returns this bigrade-$(1,1)$ polynomial; its matrix is the kernel $K$, whose
        inverse is the free propagator.

  \item[\code{fpairs} / \code{propagator\_indices}] Per-edge
        $(\text{response index},\text{physical index})$ matrix-index pairs that
        identify which field's propagator a half-edge uses. They are threaded into
        the spectral-assignment sum on the coupled multi-field path.

  \item[\code{fsolve}] SciPy's multivariate root-finder
        (\code{scipy.optimize.fsolve}), used to solve the mean-field saddle. Its
        integer convergence flag is \code{ier}; \code{ier == 1} signals clean
        convergence.

  \item[\code{fundamental} / \code{parameters}] The dict of numeric parameter
        values keyed by parameter name. \code{parameters} is the canonical name;
        \code{fundamental} is a deprecated alias that \code{Config.\_\_post\_init\_\_}
        mirrors onto the other.

\end{description}

% =====================================================================
\section*{G}
% =====================================================================
\begin{description}[leftmargin=0pt,style=nextline]

  \item[Gauss--Legendre / Gauss--Hermite / Gauss--Laguerre quadrature] Fixed-node
        numerical integration rules, each exact for polynomials up to a degree.
        \emph{Legendre} integrates over a finite interval (vertex times, the
        Schwinger $\sigma$); \emph{Hermite} integrates against a Gaussian weight
        $\ee^{-x^{2}}$ (the form-factor loop average) and an $n$-node rule is exact
        to polynomial degree $2n-1$; \emph{Laguerre} integrates
        $\int_0^{\infty}\ee^{-x}g(x)\,\dd x$ exactly for polynomial $g$ (the
        $[0,\infty)$ Schwinger parameters in the oracle path).

  \item[Generalized eigenvalue problem] The problem $C v=\lambda D v$ for matrices
        $C,D$; it reduces to the ordinary eigenproblem when $D=I$. Linear-stability
        analysis needs the general form because the ``mass'' matrix $A$ is singular
        --- algebraic constraints contribute zero rows.

  \item[\code{@dataclass}] A Python decorator that auto-generates a class's
        constructor (and a few dunder methods) from its declared fields. \code{Config}
        is a \code{@dataclass}, which is why you set only the fields you care about
        and inherit defaults for the rest.

  \item[Gillespie] The exact stochastic-simulation algorithm for jump processes (it
        schedules events one at a time). The Hawkes simulators here use a
        \emph{fixed-step Poisson-draw} approximation rather than true Gillespie
        event scheduling.

  \item[GIL (Global Interpreter Lock)] Python's lock that serialises bytecode
        execution across threads. Heavy NumPy operations \emph{release} the GIL
        while they run in C, which is precisely why thread (not fork) parallelism
        gives a real speed-up on the spatial path.

  \item[\code{G\_ft = K\_ft$^{-1}$}] The propagator matrix in frequency: the inverse
        of the quadratic-form kernel $K_{\mathrm{ft}}$. Its rows index physical
        fields and its columns index response fields; its poles are relaxation
        modes and its residues give the time-domain exponentials.

  \item[\code{graphs.trees(n)}] A Sage generator that yields every non-isomorphic
        tree on $n$ vertices exactly once. Diagram topologies are built by starting
        from these loop-free skeletons and adding edges.

  \item[GTaS] A noise/model family --- ``Gaussian-and-shifts'', a Bernoulli +
        Gaussian external-rate process modelled by \code{GTaSNoise}. It is the
        correlated-input model the cumulant machinery was originally built for; at
        order $N=2$ all its higher cumulants are local.

\end{description}

% =====================================================================
\section*{H}
% =====================================================================
\begin{description}[leftmargin=0pt,style=nextline]

  \item[Hankel transform / $J_0$] The radial Fourier transform in two dimensions,
        which uses the Bessel function $J_0$. (Spatial $d=2$ correlators.)

  \item[Hawkes process] A self-exciting point process: each spike raises the rate of
        future spikes. The simulators realise it by per-step
        $\mathrm{Poisson}(\lambda\,\dd t)$ draws fed back into the voltage.

  \item[Heat kernel] The Green's function of the diffusion operator: in $d$
        dimensions $(4\pi B t)^{-d/2}\ee^{-|x|^{2}/4Bt}$. It is the
        momentum-to-position ($q\to x$) Fourier transform of a propagator mode, and
        it is what makes the analytic inverse Fourier transform possible without a
        momentum grid.

  \item[Heat-kernel / $(k,t)$ representation] Writing propagators in
        momentum--time space, $G_R(k,t)=\theta(t)\,\ee^{-(\mu+Dk^{2})t}$. Its key
        property: each edge time acts as a Schwinger parameter, so the loop momentum
        integral becomes Gaussian.

  \item[Heaviside $\Theta(t)$] The unit step function that enforces retardation
        (a response only after its cause). The framework adopts the convention
        $\Theta(0)=0$, so the boundary $\Delta t=0$ is strictly excluded from
        retarded support.

  \item[Hermite--Genocchi formula] An identity expressing a divided difference of a
        function $f$ as a simplex integral of $f^{(n)}$. It is the bridge that turns
        the Duhamel (time-ordered) integral of the Dyson dressing into a divided
        difference. (See also Opitz theorem.)

  \item[Horner's rule] The efficient nested form for evaluating a polynomial,
        $((c_n x + c_{n-1})x + \cdots)$, using the minimum number of
        multiplications.

\end{description}

% =====================================================================
\section*{I}
% =====================================================================
\begin{description}[leftmargin=0pt,style=nextline]

  \item[Idempotent] A property of an operation: running it twice has the same effect
        as running it once. Several builder methods are idempotent so that
        re-authoring a theory does not double-apply a change.

  \item[IFT (inverse Fourier transform)] Here, the momentum-to-position transform
        $q\to x$ (or $k\to x$). ``Analytic IFT'' means it is done in closed form
        --- because each Schwinger sample's $q$-dependence is Gaussian, the
        transform is a heat kernel --- which retires the discretised $q$-grid.

  \item[Image sum / method of images] Constructing a periodic-boundary Green's
        function as a sum of displaced copies of the infinite-domain kernel.

  \item[\code{importlib.util} dynamic import] The standard-library mechanism for
        importing a Python file from an \emph{arbitrary path} rather than by a name
        on \code{sys.path}. The canonical three-step idiom ---
        \code{spec\_from\_file\_location}, \code{module\_from\_spec},
        \code{exec\_module} --- is how \code{load\_theory} imports a
        \code{.theory.py} file.

  \item[Incidence digraph (coloured)] See \emph{coloured incidence digraph}.

  \item[Interaction vertex] Any action monomial of total degree $\ge 3$ (so it has
        $\ge 1$ physical leg). It becomes a diagram vertex with $n_{\tilde{}}$
        response legs and $n_{\mathrm{phys}}$ physical legs; it is the building
        block of every loop. In code, a \code{VertexType}. Contrast a noise
        \code{SourceType} ($n_{\mathrm{phys}}=0$).

  \item[Internal vertex] A non-leaf vertex of a diagram (degree $\ge 2$); an
        interaction point. Its time is an integration variable.

  \item[\code{ipywidgets}] Jupyter's interactive-widget library, imported in the UI
        as \code{W}. It keeps Python objects on the kernel side synced to live HTML
        controls in the notebook --- sliders, text boxes, buttons.

  \item[IPython.display] IPython's rich-output API (\code{display}, \code{HTML},
        \code{Javascript}); the UI uses it to render its controls and status panel.

  \item[Isserlis / Wick theorem] The statement that the moment of a product of
        jointly Gaussian variables equals the sum over perfect matchings of products
        of pairwise covariances. It is the closed form for the loop-momentum average
        in the coupled spatial integrator.

  \item[Isomorphism (of graphs)] A relabelling of one graph's vertices that maps it
        onto another edge-for-edge. Isomorphic diagrams are ``the same diagram'';
        coloured isomorphism additionally requires colours (leaf vs internal, field
        types) to be preserved.

  \item[Iteration vs compound saddle] Two kinds of mean-field saddle. An
        \emph{iteration} (or \emph{closure}) saddle has a right-hand side that is a
        single function call (\code{nstar = phi(vstar)}) and is substituted via the
        formal rename; a \emph{compound} saddle's right-hand side is a
        parameter/saddle expression and is substituted concretely. (In legacy
        heterogeneous code, ``iteration saddle vs compound saddle'' also names the
        variables \code{fsolve} actually iterates --- rates $n^{*}$ --- versus those
        evaluated from them each step --- voltages $v^{*}$.)

  \item[\code{is\_trivial\_zero()}] Sage's \emph{syntactic} zero test: \code{True}
        only when an expression is literally zero. It is robust to the ambiguity of
        unbound parameters, where a full simplification might wrongly claim (or fail
        to claim) zero.

\end{description}

% =====================================================================
\section*{J}
% =====================================================================
\begin{description}[leftmargin=0pt,style=nextline]

  \item[$j$ (leaf surplus)] The number of ``extra'' tree leaves that the $\ell$
        added loop-edges will consume. Trees are generated with $k+j$ leaves so that
        adding $\ell$ edges (each pairing two leaves into an internal line) leaves
        exactly the $k$ external legs.

\end{description}

% =====================================================================
\section*{K}
% =====================================================================
\begin{description}[leftmargin=0pt,style=nextline]

  \item[$k$ (correlator order) / $k$-point function] The number of external legs ---
        equivalently the order of the cumulant. A $k$-point function has diagrams
        with $k$ external legs; $k=2$ is the covariance / power spectrum.

  \item[$K_0$] The modified Bessel function of the second kind of order zero; it
        appears in the static two-dimensional spatial correlator.

  \item[Keldysh / correlation \code{C} edge vs retarded \code{R} edge] A naming of
        the two edge kinds. A \code{C} (correlation) edge carries a noise-driven
        correlation, with its Schwinger parameter integrated over $[|t|,\infty)$; an
        \code{R} (retarded) edge is a fixed-duration causal propagator. (See \code{C}
        edge, \code{R} edge.)

  \item[Kernel ($g$)] A convolution (memory) kernel, e.g. a synaptic filter,
        coupling $K\star y=\int K(t-t')\,y(t')\,\dd t'$. It is declared either by its
        time form or --- preferred --- by its Fourier image $\tilde K(\omega)$.

  \item[Kernel matrix $K$] The matrix of bilinear $(1,1)$ coefficients of the free
        action, with layout $[\text{response row},\,\text{physical column}]$. Its
        inverse is the propagator $G=K^{-1}$.

  \item[Kernel normalization] The convention that synaptic / temporal kernels
        integrate to $1$. As a result, at the stationary saddle every kernel symbol
        is simply replaced by $1$.

  \item[\code{kernel\_ft\_image} hook] A model-supplied callable
        $(\mathrm{ns},\omega)\mapsto\text{subs-dict}$ that maps each abstract kernel
        symbol $g$ to its frequency image $\hat g(\omega)$. It is how a model tells
        the propagator builder what its kernels look like in Fourier space.

  \item[\code{K\_ker} (kernel form)] The kernel $K$ written with abstract
        $\delta$/$\delta'$ symbols ($c_0\,\delta + c_1\,\delta'$) so that it
        Fourier-transforms cleanly into \code{K\_ft}.

\end{description}

% =====================================================================
\section*{L}
% =====================================================================
\begin{description}[leftmargin=0pt,style=nextline]

  \item[Lag-dependent normalization] In a correlation estimator, dividing each lag's
        sum by the number of \emph{overlapping} valid bin pairs ($n_{\mathrm{valid}}
        -|\text{lag}|$), so the finite-window estimate stays unbiased.

  \item[\code{lambdify}] SymPy's utility that compiles a symbolic expression into a
        fast NumPy-callable function. (Sage's analogue is \code{fast\_callable}.)

  \item[\code{Laplacian} / \code{Lap}] The formal spatial second-derivative
        operator $\nabla^{2}$. In plain authoring it is used multiplicatively
        (\code{D*Laplacian*phi} $=D\nabla^{2}\varphi$) and is sent to $0$ on a
        spatially uniform background; in operator-IR authoring it becomes the call
        node \code{Lap()}.

  \item[Leaf / external leg] A degree-1 vertex of a diagram --- an external argument
        of the correlation function. For $k\ge 3$, one leaf (the \emph{origin leaf})
        is pinned to $t=0$.

  \item[Leg] One field ``stub'' at a vertex or source, identified by
        $(\text{field\_base\_name},\,\text{population\_index})$ and repeated once per
        exponent. Response legs feed outgoing edges; physical legs feed incoming
        edges.

  \item[Linear extension] A topological sort of a partial order --- one total
        ordering consistent with all the constraints. Each linear extension of the
        causal poset is one chain simplex / one causal chamber; the polytope is the
        disjoint union over all of them.

  \item[Linear stability] The classification of a fixed point: it is linearly stable
        when every (finite) eigenvalue of the linearised dynamics has negative real
        part, so small perturbations decay. Opt in with
        \code{.stability\_analysis(True)}; the test is the generalized eigenvalue
        problem $(B-\ii\omega A)\,\delta x=0$.

  \item[Loop number / loop order $\ell$ (\code{max\_ell})] The number of independent
        loops in a diagram, $|E|-|V|+1$ --- equivalently the order in the
        perturbative expansion. $\ell=0$ is tree level; each additional loop is one
        more power of the small parameter. \code{max\_ell} is the run's cutoff.

  \item[Lorentzian (kernel)] The exponential $\ee^{-|\tau|/\tau_c}$, named for its
        Lorentzian Fourier transform $\propto 1/(\omega^{2}+1/\tau_c^{2})$. It is
        the stationary covariance of an Ornstein--Uhlenbeck colored-noise process.

  \item[Lyapunov equation] The matrix equation $A\Sigma+\Sigma A^{\mathsf T}=Q$ (or
        $A\Sigma+\Sigma A^{\mathsf H}=Q$); its solution $\Sigma$ is the stationary
        covariance of the linear SDE $\dd x=-A x\,\dd t+\text{noise}(Q)$. Solved by
        \code{scipy.linalg.solve\_continuous\_lyapunov} (Bartels--Stewart) in the
        coupled simulators and oracle.

\end{description}

% =====================================================================
\section*{M}
% =====================================================================
\begin{description}[leftmargin=0pt,style=nextline]

  \item[\code{m} (after delta)] The number of integration variables that survive
        $\delta$-elimination within a subset. Its value selects which closed-form
        integrator runs: $m=0$, $1$, $2$, or the general $m\ge 3$
        (chain-simplex) path.

  \item[\code{manifest.json}] The human-readable index of cached entries (their
        \code{stage}, \code{k}, \code{loop\_order}, \code{saved\_at}). It is
        best-effort and \emph{not} authoritative --- the \code{.sobj} files are the
        truth.

  \item[Markovian embedding / Markovianization] Enlarging the state space with
        auxiliary fields so that a non-Markovian (memory-carrying, colored) process
        becomes Markovian (white-driven, local in time). The observable output field
        is unchanged; only the representation grows. Toggled with
        \code{.markovianize(...)}.

  \item[Mass $A$ (a.k.a. $\mu$)] The $k^{0}$ term of a spatial inverse propagator
        --- the relaxation rate. Also the reaction (mass) matrix $M$ for a
        multi-field theory: the momentum-independent part of the linearised inverse
        propagator; a diagonal $M$ means decoupled fields.

  \item[Matched-cutoff principle] When comparing a simulation against theory,
        comparing them at the \emph{same} ultraviolet cutoff $k_{\max}=\pi/\dd x$,
        because a finite-difference lattice's dispersion differs from the continuum
        $q^{2}$ at order $\dd x$.

  \item[Mean field / saddle point ($\varphi^{*}$)] The deterministic steady state
        about which the loop expansion is taken --- the constant background that
        makes the action's field-linear (tadpole) term vanish. Saddle parameters end
        in \code{star} (e.g. \code{xstar}, \code{nstar}); they are solved by the
        engine, never supplied. The saddle (MF) sector is the bigrades
        $(0,0),(1,0),(0,1)$, which the framework verifies and zeroes.

  \item[Merged-residue tensor $B_\alpha$] The grouped path's per-pole-tuple
        coefficient, $\sum_j cp_j\prod\delta\prod C^{(j)}$ --- a single tensor that
        replaces the per-diagram Cartesian product $\prod_e C_{\alpha_e}$. Summing
        residues once, across diagrams sharing a topology, is what makes the grouped
        Phase~J fast and exact.

  \item[Mode $(C_\alpha,\lambda_\alpha)$ / mode sum] One pole's residue $C_\alpha$
        and rate $\lambda_\alpha=\ii p_\alpha$, so that the smooth part of a
        propagator is $\sum_\alpha C_\alpha\ee^{\lambda_\alpha\Delta t}$. This
        per-edge representation is the \emph{mode sum}, stored as an
        \code{EdgeModeSum}.

  \item[Moment] A raw correlation $\langle\varphi\cdots\varphi\rangle$, or the same
        for the centred field (a central moment). Reconstructed from cumulants by the
        set-partition (cluster-expansion) formula.

  \item[\code{mp.dps}] mpmath's decimal-digit precision setting (raised to $50$ in
        the precision-rescue paths). (See mpmath.)

  \item[mpmath] A Python library for arbitrary-precision floating point and special
        functions. It is used where machine \code{double} would suffer catastrophic
        cancellation --- e.g. the \code{erf}-cancellation in the heat-kernel closed
        form.

  \item[\msrjd{} action / field theory] The
        Martin--Siggia--Rose--Janssen--De~Dominicis path-integral representation of a
        stochastic dynamical system. (See the boxed definition below.)

\end{description}

\begin{defn}
The \msrjd{} construction rewrites a Langevin (stochastic differential)
equation as a path integral with weight $\ee^{-S[\varphi,\tilde\varphi]}$. Every
\textbf{physical field} $\varphi$ is paired with a conjugate \textbf{response
field} $\tilde\varphi$ (the ``tilde'' or ``hatted'' field; written with a
\code{t}- or \code{d}-prefix in code, e.g.\ \code{xt}, \code{dh}). The
deterministic drift sits in the $\tilde\varphi\cdot(\dots)$ terms; the noise sits
in the response-field-only terms (white Gaussian noise is the single vertex
$-D\,\tilde\varphi^{2}$, the bigrade-$(2,0)$ sector). Sourcing the response field
produces response functions; sourcing the physical field produces correlation
functions. Propagators are directed (causal): the retarded $\Gret$ connects a
physical leg to a response leg and is non-zero only for $t>0$.
\end{defn}

\begin{description}[leftmargin=0pt,style=nextline]
  \item[Multi-edge (multigraph)] Two or more edges between the same pair of
        vertices; a double edge forms a ``bubble'' loop. Enabled in Sage with
        \code{multiedges=True}.

  \item[Multi-start Newton] Running a Newton-type root finder from many random
        initial guesses to discover \emph{all} the roots of a nonlinear system (e.g.
        both wells of a double-well saddle), then sorting them so
        \code{fixed\_point\_index} can pick one deterministically.

  \item[\code{multiset\_permutations} (SymPy)] A SymPy generator yielding the
        $N!/\prod_r n_r!$ distinct orderings of a multiset (a list with repeats). It
        enumerates the distinct external-leaf assignments without double-counting
        identical legs.
\end{description}

% =====================================================================
\section*{N}
% =====================================================================
\begin{description}[leftmargin=0pt,style=nextline]

  \item[\code{nauty}] Brendan McKay's C library for graph canonical labelling and
        isomorphism testing --- the name stands for ``No AUTomorphisms, Yes?''. It is
        bundled inside SageMath and is invoked implicitly by every
        \code{is\_isomorphic}, \code{canonical\_label}, and \code{automorphism\_group}
        call. It is what makes diagram deduplication and symmetry-factor computation
        tractable.

  \item[\code{networkx}] A pure-Python graph library used inside the engine for some
        graph manipulations (distinct from Sage's \code{Graph}/\code{DiGraph}).

  \item[Newton refinement] Polishing an approximate root of $\det(K(\omega))$ to
        machine precision with Newton's method, after a coarser method has located
        it.

  \item[Noise kernel] The response-field-only sectors of the action --- bigrades
        $(n_{\tilde{}}\ge 2,\,0)$. These are the noise sources from which the diagram
        expansion's stochastic forcing comes.

  \item[NoiseSourceType / non-local cumulant kernel] A noise vertex whose kernel
        $\kappa(\tau)$ introduces \emph{extra} $\tau$ integration variables and a
        non-rational integrand, which disables the analytic (closed-form) integration
        path and forces \code{scipy.nquad}. Contrast a plain (local, white) noise
        source.

  \item[\code{nquad}] SciPy's nested numerical multidimensional quadrature
        (\code{scipy.integrate.nquad}); the slow adaptive fallback (about
        $10^{-8}$ precision) that the analytic machinery exists to avoid.

  \item[\code{ns} (NamespaceForFields)] The runtime object handed to the model's
        lambdas, carrying the symbolic field and parameter variables (Sage \code{SR}
        objects) as attributes. It is built by the consumer, so a model author can
        write \code{ns.mu}, \code{ns.x}, and so on.

  \item[\code{num\_params}] The substitution dictionary
        $\{\code{SR}\ \text{symbol}\to\text{float}\}$ that converts every symbolic
        propagator, vertex, and residue into a numeric one. It is the deliverable of
        the mean-field / numericisation stage. Its keys are Sage \code{SR} symbols,
        not strings.

  \item[\code{numba} / \code{@njit} / nopython mode] \code{numba} is a
        just-in-time compiler that turns annotated Python functions into native
        machine code. Its \code{@njit} decorator forces ``nopython'' mode --- no
        interpreter fallback --- giving roughly $100\times$ speed-ups on tight
        numeric loops (e.g. the chain-simplex inner loop, the Euler--Maruyama
        stepper). It cannot type a Sage ring element, which is why simulator inputs
        must be pre-cast to plain Python floats.

  \item[NumPy \code{ndarray}] The $n$-dimensional array type underlying all numeric
        data in the pipeline --- grids, correlators, residue matrices.

\end{description}

% =====================================================================
\section*{O}
% =====================================================================
\begin{description}[leftmargin=0pt,style=nextline]

  \item[\code{.operator()} / \code{.operands()}] Sage methods that return,
        respectively, an expression's outermost operation and the list of its
        arguments. Walking an expression with these is how the framework inspects a
        symbolic tree without string parsing.

  \item[Operator IR] The intermediate representation for derivative interaction
        vertices: $\nabla^{2}$ (\code{Lap}), $\partial_t$ (\code{Dt}), and
        $\partial_{x_i}$ (\code{Dx}) become inert symbolic \emph{call} nodes with
        explicit algebra passes (\file{pipeline/spatial\_operator\_ir.py}). It is
        what lets spatial theories with derivative vertices (Model~B, KPZ, Burgers)
        carry the right momentum form factors. Authoring is gated behind
        \code{.operator\_ir()}.

  \item[Opitz theorem] The fact that a divided difference equals the top-right entry
        of $f$ applied to the upper-bidiagonal ``Opitz matrix'' (the nodes on the
        diagonal, ones on the superdiagonal). It is \emph{confluent-safe} --- no
        $0/0$ at repeated nodes --- which is why the Dyson dressing uses it for the
        $\Phi_n$ primitive.

  \item[ORACLE-ONLY] A label on a module that exists solely as an \emph{independent}
        numerical cross-check, kept off the production path. The codebase
        banner-annotates such modules.

  \item[Orbit--stabilizer theorem] The group-theory identity
        $|\text{orbit}|=|\text{group}|/|\text{stabilizer}|$. It justifies the
        external-Wick compensation: the number of distinct external-leaf labellings
        is the automorphism group's size divided by the stabilizer that fixes the
        leaves.

  \item[\code{origin\_leaf\_idx} / origin pinning] The index of the external leaf
        pinned to $t=0$. By time-translation invariance a $k$-point cumulant depends
        only on time differences, so one leaf is fixed at the origin and the
        resulting callable computes a function of the remaining differences.

  \item[Ornstein--Uhlenbeck (OU) process] The simplest mean-reverting linear SDE,
        $\dd\xi/\dd t=-\xi/\tau_c+\text{white noise}$. It is Gaussian, with
        stationary variance $D/\mu$ and a single-Lorentzian (exponential)
        autocorrelation. It appears both as a \emph{theory} (the linear baseline)
        and as the auxiliary \emph{colored-noise generator} in Markovian embedding.

\end{description}

% =====================================================================
\section*{P}
% =====================================================================
\begin{description}[leftmargin=0pt,style=nextline]

  \item[Partition (\code{nauty}/Sage sense)] A vertex colouring passed to the graph
        algorithms; vertices may only be mapped within the same colour class. (Not to
        be confused with a set partition of indices.)

  \item[Per-leg vs composite] See \emph{composite vs per-leg vertex}.

  \item[Phase J] The time-domain loop-integration stage of the temporal pipeline
        (\code{compute\_correction\_td}). (See the boxed definition below.)

  \item[$\Phi_n(t;\nu)$] The confluent-safe nested-simplex time primitive of the
        Dyson dressing, with $\Phi_0=1$ and $\Phi_1=(1-\ee^{-t\nu})/\nu$. It is
        evaluated via the Opitz matrix to stay stable at repeated nodes.

  \item[$\mathcal{H}_n(t)$] The matrix factor of the $n$-th Dyson order: a sum over
        projector strings times $\ee^{-m_{\alpha_0}t}\,\Phi_n$.

  \item[$\varphi^{(k)}$ / \code{phik\_i}] The $k$-th derivative of the transfer
        function at the saddle of population $i$; \code{phi0\_i}
        $=\varphi(v^{*}_i)=n^{*}_i$. These derivatives become the diagram couplings.

  \item[\code{\_PhiCallableList}] A list of per-population transfer-function
        callables that also supports a bare scalar call \code{f(v)} when there is
        exactly one population.

  \item[Physical (fluctuation) field] The actual fluctuating quantity
        $\delta\varphi=\varphi-\varphi^{*}$, the ``in-leg'' of a diagram. Natural code
        names: \code{n}, \code{v}, \code{x}, \code{h}. (See also fluctuation field.)

  \item[Picklable] Serializable by Python's \code{pickle}, so that
        \code{multiprocessing} can ship an object between processes. This is why the
        CGF kernel callables (\code{\_CGFKernelCallable}, \code{\_CGFKernelSum}) are
        defined as module-level classes rather than closures.

  \item[PipelineCache] The low-level cache store keyed by \code{(stage, k, ell)},
        writing \code{.sobj} files with a JSON manifest
        (\file{msrjd/core/cache.py}).

  \item[Pole $\omega_k$ / pole-residue form] A retarded zero of
        $\det(K_{\mathrm{ft}})$ with $\operatorname{Im}(\omega_k)>0$, setting a
        relaxation/oscillation rate. The propagator's \emph{pole-residue form} is
        $\sum_k C_k\,\ee^{-\ii\omega_k t}+D_\delta\,\delta(t)$ --- the form Phase~J
        integrates against. In code, \code{pole\_vals} holds the poles and
        \code{C\_mats} the matrix residues.

  \item[Pole tuple $\alpha$] A choice of one pole per smooth edge. Expanding the
        product of per-edge mode sums yields a sum over all pole tuples; each tuple's
        contribution is a single exponential whose polytope integral is closed-form.

  \item[Polytope / chamber] The convex region carved out by the causal
        $\Delta t>0$ constraints --- the integration domain. (See causal poset, chain
        simplex, causal chamber.)

  \item[\code{poly.dict()}] A Sage polynomial-ring method returning
        $\{\text{exponent tuple}\to\text{coefficient}\}$ --- the pickle-stable
        ``dict form'' the expand cache stores in place of the (less stable)
        polynomial objects themselves.

  \item[\code{PolynomialRing(SR, names)} / fraction field] A polynomial ring over the
        Symbolic Ring whose generators are the named field symbols; it exposes
        monomials via \code{.dict()}. Its fraction field $\mathrm{Frac}(R[\omega])$
        reduces each propagator entry to a canonical irreducible $P/Q$.

  \item[Population] A named index set --- a group of \code{size N} identical units.
        Attaching a field to a population makes it a length-$N$ vector and introduces
        \code{for i in pop} sums in the action. Heterogeneous theories carry several
        populations.

  \item[\code{pop\_idx} (population index)] The 1-based subscript on a field
        generator (e.g. the \code{2} in \code{nt2}), distinguishing
        populations/components.

  \item[Precompute] The cheap one-time structural pass --- \code{expand} at order 2,
        the saddle sanity check, the mean-field solve, and the propagator build ---
        that populates the durable caches. Exposed as \code{pipeline.precompute(model)}
        and run by the UI's \emph{Pre-compute} button.

  \item[Prediagram] A bare, oriented graph \emph{topology} --- leaves, internal
        vertices, and arrows on edges --- carried as the tuple
        \code{(D, G, leaves, internal)}, \emph{before} any propagator types or vertex
        monomials are assigned. It is the output of the enumeration subsystem.

  \item[Preparse / SageMath preparser] Sage's source rewrite of notebook-cell text:
        integer literals become \code{Integer}, floats become \code{RealNumber}, and
        \code{\^{}} becomes \code{**}. Because this rewrite changes plain Python
        semantics, the simulators must live in (un-preparsed) \code{.py} files and
        accept pre-cast scalars.

  \item[Proper / strictly proper rational] A rational function with
        $\deg(\text{numerator})<\deg(\text{denominator})$; it vanishes at infinity
        (no $\delta$ term). A non-proper rational has a polynomial part, which becomes
        a $\delta(t)$ in the time domain.

  \item[Propagator $G=K^{-1}$] The two-point Gaussian correlation/response function
        --- the building block of every Feynman edge --- obtained as the matrix
        inverse of the bilinear kernel $K$. (See \code{G\_ft}, retarded propagator.)

\end{description}

\begin{defn}
\textbf{Phase J} is the temporal pipeline's loop-integration stage,
\code{compute\_correction\_td}. It performs the actual vertex-time quadrature for
each diagram, in the time domain, using the propagator's pole-residue form
$\sum_\alpha C_\alpha\ee^{\ii p_\alpha t}+D_\delta\,\delta(t)$. The product of
retarded propagators over the edges is expanded into $2^{|E|}$ choices of
``instantaneous $\delta$ vs smooth'' per edge; each $\delta$-edge eliminates one
integration variable, leaving $m$ variables that select a closed-form path
($m=0,1,2$) or the general chain-simplex path ($m\ge 3$). The \emph{grouped}
variant (Stage~4a) sums residues once across diagrams sharing a topology, via the
merged-residue tensor $B_\alpha$, dropping the precision floor to machine
$\varepsilon$.
\end{defn}

% =====================================================================
\section*{Q}
% =====================================================================
\begin{description}[leftmargin=0pt,style=nextline]

  \item[\code{q}-grid] A discretised momentum axis used by a \emph{numerical} Fourier
        transform. It is ``retired'' by the analytic heat-kernel inverse Fourier
        transform for the plain-vertex and $d=1$ cases, which removes the associated
        cutoff and ringing artefacts.

  \item[\code{QQ}] Sage's notation for the field of rational numbers $\mathbb{Q}$;
        used (alongside \code{CyclotomicField(4)}) for exact symbolic arithmetic.

  \item[\code{quad}] SciPy's adaptive one-dimensional quadrature
        (\code{scipy.integrate.quad}); the 1-D counterpart of \code{nquad}.

  \item[$Q_{\mathrm{eff}}$] The second Symanzik polynomial in its Schur-complement
        form, $Q-N^{\mathsf T}\Lambda^{-1}N$ --- the residual quadratic form carrying
        external-momentum dependence after the loop block is integrated out. (See
        Symanzik polynomials, Schur complement.)

\end{description}

% =====================================================================
\section*{R}
% =====================================================================
\begin{description}[leftmargin=0pt,style=nextline]

  \item[\code{R} edge / retarded line] A $\Gret$ propagator segment, carrying a
        causal $\theta(w\ge 0)$. (Contrast the \code{C} correlation edge.)

  \item[Reaction (mass) matrix $M$] See \emph{mass}.

  \item[Reference diffusion $D_0$] The scalar part split off the diffusion matrix,
        $\mathcal{D}=D_0 I+\hat{\mathcal{D}}$ (default $D_0=\operatorname{trace}/N$),
        chosen to make the residual $\hat{\mathcal{D}}$ small so the Dyson series
        converges quickly.

  \item[Regime 1 (finite cutoff)] Evaluating the spatial correlator \emph{at} a
        finite momentum cutoff $k_{\max}$; the continuum limit is
        $k_{\max}\to\infty$.

  \item[Residual diffusion $\hat{\mathcal{D}}$] $\mathcal{D}-D_0 I$, the Dyson
        perturbation. If it is zero, no series is needed and the tree propagator is
        exact.

  \item[Residue matrix \code{C\_mats[k]}] $\ii\cdot\operatorname{Res}_{\omega_k}
        G_{\mathrm{ft}}$ --- the amplitude matrix of mode $k$. Together with the poles
        \code{pole\_vals} it gives the time-domain propagator
        $\sum_k C_k\,\ee^{\ii\omega_k t}$.

  \item[Response field (tilde / hat field)] The conjugate field $\tilde\varphi$ of
        the \msrjd{} construction --- the ``out-leg'' in diagram language. Its code
        base names carry a \code{t}- or \code{d}-prefix (\code{xt}, \code{nt},
        \code{dv}, \code{dn}). A degree-2 monomial in it is white Gaussian noise; a
        degree-$\ge 3$ monomial is non-Gaussian noise.

  \item[Retarded propagator $\Gret$] The causal Green's function:
        $\Gret(t)=D_\delta\,\delta(t)+\Theta(t)\sum_k C_k\,\ee^{\ii\omega_k t}$. It is
        the linear response of a physical field to a response-field source ---
        non-zero only for $t>0$ (effect after cause). In momentum--time space,
        $\Gret(k,t)=\theta(t)\,\ee^{-(\mu+Dk^{2})t}$; for a multi-field theory,
        $\Theta(t)\,\ee^{-(M+\mathcal{D}k^{2})t}$.

  \item[Routing coefficients $(a_e,b_e)$] The integer ($\pm 1$ or $0$) coefficients
        in $k_e=\sum a_{ei}\ell_i+\sum b_{ej}q_j$ --- how the loop momenta $\ell$ and
        external momenta $q$ thread each edge. They are found by solving momentum
        conservation (e.g. with SymPy \code{linsolve}).

\end{description}

% =====================================================================
\section*{S}
% =====================================================================
\begin{description}[leftmargin=0pt,style=nextline]

  \item[Saddle / mean-field parameter] See \emph{mean field}. A \code{<field>star}
        parameter holding a field's steady-state value, auto-created with each
        \code{physical\_field} and solved (never passed in).

  \item[SageMath / \code{SR} (Symbolic Ring)] SageMath is an open-source mathematics
        system that bundles many libraries (including \code{nauty}, Singular, Maxima,
        Pynac/GiNaC) behind a single Python API. \code{SR} is its Symbolic Ring ---
        the universal type of symbolic expressions; \code{SR.var('x')} makes a
        symbolic variable, and \code{diff}, \code{.subs}, \code{.simplify\_full}, and
        \code{float} operate on \code{SR} elements. The physics-facing algebra of the
        pipeline lives in \code{SR}, and a Sage kernel is required to run it.

  \item[\code{sage\_eval}] Sage's namespaced string-to-symbolic-expression evaluator
        --- it parses a string into an \code{SR} expression within a supplied
        namespace of allowed symbols.

  \item[\code{sage\_solve}] Sage's symbolic equation solver (backed by Maxima); used,
        for instance, to solve the linear $\Delta t=0$ equations of
        $\delta$-elimination.

  \item[Schur complement] For a block matrix, the residual form
        $Q-N^{\mathsf T}\Lambda^{-1}N$ left after integrating out the $\Lambda$
        (loop-momentum) block --- which is exactly the second Symanzik polynomial
        $Q_{\mathrm{eff}}$.

  \item[Schwinger parameter] An auxiliary integration variable (here an edge time
        $w$, or a noise-source time $\sigma\in[0,\infty)$) introduced to turn a
        propagator denominator into an exponential. Doing this for every edge makes
        the loop \emph{momentum} integral Gaussian, which then collapses to the
        Symanzik polynomials. In the heat-kernel representation the Schwinger
        parameter \emph{is} the edge's time duration.

  \item[Self-consistency equation] The fixed-point condition the saddle must satisfy
        (e.g. $n^{*}=\varphi(E+w\cdot n^{*})$); solving it is the mean-field solve.

  \item[Self-energy $\Sigma$] The one-particle-irreducible (1PI) two-point
        correction; it dresses the propagator through the Dyson equation. In the
        spatial diagnostics, \code{Sigma} is the one-loop self-energy used for the
        tadpole mass-shift.

  \item[Self-loop] An edge from a vertex to itself; disabled in the enumeration
        (\code{loops=False}) because it is never physical here.

  \item[Set partition / cluster expansion] The combinatorial identity expressing
        moments as sums over partitions of products of cumulants over the blocks. It
        is how \code{output='moment'} reconstructs full moments from the engine's
        native cumulants.

  \item[Shot noise / contact / diagonal term] The $1/\dd t_{\mathrm{bin}}$-scaling
        $\delta$-function piece of a point-process correlator --- a spike correlated
        with itself. It is removed by the factorial-moment correction. The temporal
        analogue is a \emph{shot-noise $\delta$-spike}: a residual same-population
        equality $\tau=0$ left after $\delta$-elimination.

  \item[\code{.simplify\_full()}] Sage's most aggressive symbolic simplifier. It is
        powerful but slow, and it can stall on unbound parameters --- hence the
        \code{cysignals} time-budget around it and the cheaper
        \code{is\_trivial\_zero()} where a full simplify is overkill.

  \item[Slice / full grid] For $k\ge 3$, a \emph{slice} is an axis-parallel 1-D cut
        (sweep one leg's lag, hold the rest) while the \emph{full grid} is the entire
        $(k-1)$-dimensional tensor of $(x_j,\tau_j)$ events.

  \item[Slug / slugify] A filesystem-safe lowercase identifier derived from a theory
        name, \code{re.sub(r'[\^{}A-Za-z0-9]+','\_',name).strip('\_').lower()}
        (\file{pipeline/theory\_serialize.py:96}). It is the per-theory cache
        directory name and the \code{<slug>.theory.py} filename.

  \item[\code{.sobj}] ``Sage Object'': a gzip-compressed pickle written by
        \code{sage.all.save} and read by \code{sage.all.load}. It is the on-disk
        format for every cached artefact --- the expanded action, the propagator, the
        enumerated diagrams.

  \item[Source / source vertex] An internal vertex with in-degree $0$ (and not a
        leaf) --- a pure response-field monomial from the noise kernel that
        \emph{injects} noise (a \code{SourceType}, with $n_{\mathrm{phys}}=0$).

  \item[\code{spatial\_grid} / \code{spatial\_points}] How a spatial run specifies its
        evaluation points. For $k=2$, a grid of separations $x$; for $k\ge 3$, an
        explicit \code{(n\_pts, k-1, 2)} array of $(x_j,\tau_j)$ events.

  \item[\code{spatial\_info} / modes / \code{Sigma}] The spatial driver's diagnostic
        dict. \code{modes} are $(A,B,N)$ triples parameterising each propagator
        mode's mass $A$, dispersion $B$, and noise weight $N$; \code{Sigma} is the
        one-loop self-energy.

  \item[Spectral assignment / spectral projector $P_\alpha$] On the coupled
        (unequal-$D$) path, a labelling of every propagator segment with one of the
        $N_{\mathrm{modes}}$ spectral eigenmodes of the mass matrix $M=\sum_\alpha
        m_\alpha P_\alpha$; the diagram sums over all such labellings. $P_\alpha$ is
        the (rank-1 or higher) matrix projecting onto the $\alpha$-th eigenspace,
        built from the eigenvectors.

  \item[Spec dict] The flat Python dict the UI produces from its form (via
        \code{\_collect}) and the serializer renders to / parses from a
        \code{.theory.py}.

  \item[Specialization] A model-supplied substitution that replaces an abstract
        derivative symbol with a concrete value, e.g. \code{phi2\_i} $\to 2a$.

  \item[Structure factor $S(q)$] The equal-time momentum-space two-point function
        $\langle|\varphi_q|^{2}\rangle$ of a spatial field; at tree level
        $S(q)\to T/(\mu+Dq^{2})$.

  \item[Sunset] The canonical two-loop ($L=2$) test diagram: three lines run between
        two vertices.

  \item[Superset theorem] The fact that a Taylor expansion at order $N$ contains, on
        the degrees it shares, a byte-identical copy of the expansion at any lower
        order $M\le N$ (\code{by\_tp@N} $\supseteq$ \code{by\_tp@M}). It is the
        justification for the expand cache and for downgrade-filtering.

  \item[Sutherland--Hodgman] The polygon-clipping algorithm used to build the
        two-dimensional ($m=2$) integration region.

  \item[Sweep axis] In the cumulant estimator, the single external leg whose lag is
        left as \code{None}; the cumulant is reported as a function of \emph{its} lag,
        with the other legs held fixed.

  \item[Symanzik polynomials ($\Usym$, $\Fsym$/$Q_{\mathrm{eff}}$)] The two graph
        polynomials a Feynman loop integral reduces to once Schwinger/Feynman
        parameters are introduced. (See the boxed definition below.)

  \item[Symmetry factor $\Sym(\Gamma)$] The Feynman-rule combinatorial weight of a
        diagram. (See the boxed definition below.)

  \item[\code{sympy}] A lighter, pure-Python computer-algebra system (distinct from
        Sage's \code{SR}). The engine uses it for targeted jobs --- \code{lambdify} to
        compile expressions to fast NumPy functions, \code{linsolve} to route
        momenta, \code{multiset\_permutations} to enumerate leaf orderings.

  \item[\code{\_\_slots\_\_}] A Python class optimisation that removes the per-instance
        \code{\_\_dict\_\_} (saving memory and speeding attribute access). A class with
        \code{\_\_slots\_\_} needs explicit
        \code{\_\_getstate\_\_}/\code{\_\_setstate\_\_} to be picklable.

  \item[\code{substitute\_function}] A Sage method that replaces every occurrence of a
        formal function with the result of a Python callback. It is how the formal
        convolution \code{Conv} is reduced to a concrete expression.

  \item[Squarefree polynomial] A polynomial with no repeated roots. A non-squarefree
        characteristic denominator $Q$ signals a multiplicity-$\ge 2$ (Jordan-block)
        pole, which the single-pole residue formula cannot handle.

\end{description}

\begin{defn}
The \textbf{symmetry factor} $\Sym(\Gamma)$ is the combinatorial weight a diagram
$\Gamma$ carries under the Feynman rules,
\[
  \Sym(\Gamma)\;=\;\frac{\prod_{\text{legs}} n_{\text{leg}}!}{\bigl|\Aut_{\text{fixed-ext}}(\Gamma)\bigr|},
\]
the product of per-vertex Wick leg-factorials divided by the order of the
automorphism group of the coloured incidence digraph with external leaves pinned.
In the code it is carried by \code{combinatorial\_factor} (and \code{pv} in the
spatial integrator is $\Sym(\Gamma)$ times the vertex-coupling prefactor). Under
the framework's ``Path~A'' convention the dedup multiplicity is \emph{not}
multiplied back in --- $\Sym(\Gamma)$ already accounts for the size of the
isomorphism class. The automorphism orders are computed by \code{nauty} through
Sage's \code{automorphism\_group}.
\end{defn}

\begin{defn}
The \textbf{Symanzik polynomials} arise when a Feynman loop integral is rewritten
using one Schwinger parameter per edge. The first, $\Usym=\det\Lambda$, is a sum
over the diagram's spanning trees and is purely topological; the second, $\Fsym$
(equivalently $Q_{\mathrm{eff}}$, the Schur complement $Q-N^{\mathsf
T}\Lambda^{-1}N$), carries the external-momentum dependence. The Gaussian
integral over the loop momenta collapses \emph{exactly} to a ratio built from
$\Usym$ and $\Fsym$, turning a multidimensional momentum integral into an
integral over the bounded Schwinger parameters. The degenerate direction
$\det\Lambda=\Usym\to 0$ (e.g. a self-loop) is the $\Usym^{-d/2}$ singularity
that would make naive Monte-Carlo sampling infinite-variance; \Daedalus{} avoids
it by an analytic radial-Bessel reduction.
\end{defn}

% =====================================================================
\section*{T}
% =====================================================================
\begin{description}[leftmargin=0pt,style=nextline]

  \item[Tadpole] Two related senses. (1) The linear-in-fluctuation term of the
        expanded action --- the bigrade-$(1,0)$ piece --- whose vanishing \emph{is}
        the saddle equation. (2) In diagrams, a $C$ self-loop ($u=v$): a decoupled
        $\langle\varphi^{2}\rangle$-type loop that carries no external momentum.

  \item[$\tau_c$ (\code{tauc})] The noise correlation time, equivalently the
        Ornstein--Uhlenbeck relaxation time of the colored-noise generator.

  \item[$\tau_j=t_j-t_0$] The $k-1$ time differences a $k$-point cumulant depends on
        (by time-translation invariance).

  \item[Taylor order (\code{taylor\_order}, $N$)] The total-degree truncation of the
        multivariate Taylor expansion of the action in the fluctuation fields. It
        equals the maximum number of legs a vertex can have, and defaults to
        $\max(k+2\,\code{max\_ell},2)$. It must be at least the largest prediagram
        vertex degree --- the job of \code{degree\_scan.ensure\_taylor\_order}.

  \item[\code{ThreadPoolExecutor} / \code{as\_completed}] Python standard-library
        thread-pool primitives used for opt-in parallelism --- threads, never fork,
        on macOS-safety grounds. \code{as\_completed} yields futures as they finish.

  \item[Tier 1 vs Tier 2] In the spatial pipeline, Tier~1 is the diagonal
        closed-form heat-kernel path; Tier~2 is the numerical / coupled /
        higher-derivative fallback (the v2 integrator).

  \item[Topology] An \emph{undirected} prediagram --- vertices and edges but no
        arrows yet. It is the intermediate output of the topology stage, before
        orientation makes it a prediagram.

  \item[Transfer function ($\varphi$)] The (possibly nonlinear) firing-rate map from
        input (voltage/current) to output (rate); in code, the user function literally
        named \code{phi}. Its Taylor coefficients at the saddle,
        \code{phi0\_i, phi1\_i, ...}, become the vertex couplings. (The name clashes
        with the field $\varphi$; context disambiguates.)

  \item[Tree] A connected acyclic graph; the loop-free skeleton from which
        topologies are built by adding edges.

  \item[Trust chain] The closed validation loop the repository runs on: an analytic
        diagram prediction versus an \emph{independent} brute-force simulation.
        Agreement within error bars is what licenses every analytic result.

  \item[Typed diagram (\code{TypedDiagram})] A prediagram fully decorated: a definite
        field type on every leg, a propagator $(\text{resp idx},\text{phys idx})$ on
        every edge, and a coupling on every vertex. Many typed diagrams can share one
        prediagram.

\end{description}

% =====================================================================
\section*{U}
% =====================================================================
\begin{description}[leftmargin=0pt,style=nextline]

  \item[$\Usym$ (first Symanzik polynomial)] $\det\Lambda$, the sum over the graph's
        spanning trees. (See the Symanzik boxed definition under \textbf{S}.) Its
        vanishing, $\det\Lambda\to 0$, is the source of the $\Usym^{-d/2}$
        Monte-Carlo singularity, cured analytically by the radial-Bessel reduction.

\end{description}

% =====================================================================
\section*{V}
% =====================================================================
\begin{description}[leftmargin=0pt,style=nextline]

  \item[Vertex] An interaction term of the action (cubic and higher in the fields)
        that becomes a diagram vertex. (See interaction vertex.)

  \item[Vertex coloring / colored isomorphism] Assigning colours to vertices (here
        $0=$ leaf, $1=$ internal) and requiring an isomorphism to \emph{preserve}
        them, so a leaf can only map to a leaf. Implemented via Sage's
        \code{set\_vertex}.

  \item[\code{vertex\_form\_factor\_signature}] See \emph{form-factor signature}
        (\file{pipeline/\_expand\_cache.py:208}).

  \item[\code{vertex\_role\_signature}] A depth-1 graph invariant of a vertex
        (\file{msrjd/diagrams/symmetry.py:247}); it was the basis of the
        \emph{old}, now-superseded external-Wick compensation heuristic.

  \item[\code{v2\_max} / \code{v3\_max}] Search bounds on the number of degree-2 and
        degree-$\ge 3$ vertices a topology may have: $\code{v3\_max}=k+j-2$,
        $\code{v2\_max}=k+3\ell-j-1$. (See completeness bound.)

\end{description}

% =====================================================================
\section*{W}
% =====================================================================
\begin{description}[leftmargin=0pt,style=nextline]

  \item[White noise] Noise with zero correlation time; its covariance is
        $\propto\delta(\tau)$. In the \msrjd{} action it is the single local vertex
        $-D\,\tilde\varphi^{2}$ (the bigrade-$(2,0)$ sector). Contrast colored noise.

  \item[Wick theorem] See \emph{Isserlis / Wick theorem}.

\end{description}

% =====================================================================
\section*{Z}
% =====================================================================
\begin{description}[leftmargin=0pt,style=nextline]

  \item[$\Theta(0)=0$ convention] The convention that the retarded support boundary
        $\Delta t=0$ is strictly excluded, enforced by strict inequalities and the
        Heaviside filter. (Filed under Z for ``zero''; see Heaviside.)

  \item[\code{ZMQInteractiveShell}] The IPython shell class used only by a
        ZMQ/Jupyter kernel. The fork-safety guard fingerprints a notebook by
        detecting this class --- if it is the active shell, the fork-based parallel
        path is downgraded to serial (because forking a Jupyter kernel after
        Cocoa/BLAS init crashes it).

  \item[$2^{-n_C}$ (convention factor)] The factor converting the enumeration's
        $2T$ noise-vertex normalisation to the spatial integrator's unit-amplitude
        $C$ edges, with $n_C$ the number of correlation edges.

\end{description}
